%% For double-blind review submission, w/o CCS and ACM Reference (max submission space)
\documentclass[acmsmall,review,anonymous,dvipsnames]{acmart}\settopmatter{printfolios=true,printccs=false,printacmref=false}
%% For double-blind review submission, w/ CCS and ACM Reference
%\documentclass[acmsmall,review,anonymous]{acmart}\settopmatter{printfolios=true}
%% For single-blind review submission, w/o CCS and ACM Reference (max submission space)
%\documentclass[acmsmall,review]{acmart}\settopmatter{printfolios=true,printccs=false,printacmref=false}
%% For single-blind review submission, w/ CCS and ACM Reference
%\documentclass[acmsmall,review]{acmart}\settopmatter{printfolios=true}
%% For final camera-ready submission, w/ required CCS and ACM Reference
%\documentclass[acmsmall]{acmart}\settopmatter{}

\RequirePackage[utf8]{inputenc}

%% Journal information
%% Supplied to authors by publisher for camera-ready submission;
%% use defaults for review submission.
\acmJournal{PACMPL}
\acmVolume{1}
\acmNumber{CONF} % CONF = POPL or ICFP or OOPSLA
\acmArticle{1}
\acmYear{2018}
\acmMonth{1}
\acmDOI{} % \acmDOI{10.1145/nnnnnnn.nnnnnnn}
\startPage{1}

%% Copyright information
%% Supplied to authors (based on authors' rights management selection;
%% see authors.acm.org) by publisher for camera-ready submission;
%% use 'none' for review submission.
\setcopyright{none}
%\setcopyright{acmcopyright}
%\setcopyright{acmlicensed}
%\setcopyright{rightsretained}
%\copyrightyear{2018}           %% If different from \acmYear

%% Bibliography style
\bibliographystyle{ACM-Reference-Format}
\DeclareUnicodeCharacter{1D706}{$\lambda$}
\DeclareUnicodeCharacter{209B}{${}_s$}
%% Citation style
%% Note: author/year citations are required for papers published as an
%% issue of PACMPL.
\citestyle{acmauthoryear}   %% For author/year citations


%%%%%%%%%%%%%%%%%%%%%%%%%%%%%%%%%%%%%%%%%%%%%%%%%%%%%%%%%%%%%%%%%%%%%%
%% Note: Authors migrating a paper from PACMPL format to traditional
%% SIGPLAN proceedings format must update the '\documentclass' and
%% topmatter commands above; see 'acmart-sigplanproc-template.tex'.
%%%%%%%%%%%%%%%%%%%%%%%%%%%%%%%%%%%%%%%%%%%%%%%%%%%%%%%%%%%%%%%%%%%%%%


%% Some recommended packages.
\usepackage{booktabs}   %% For formal tables:
                        %% http://ctan.org/pkg/booktabs
\usepackage{subcaption} %% For complex figures with subfigures/subcaptions
                        %% http://ctan.org/pkg/subcaption

\usepackage{mathpartir}

\usepackage{stmaryrd}
\usepackage{colortbl}

\usepackage{bussproofs}
%%\usepackage{semantic}
\usepackage{stmaryrd}
\usepackage{mathtools}
\usepackage{amsmath}
\usepackage{verbatim}
\usepackage[capitalize]{cleveref}   % \Cref

% Boxed figures
\usepackage{float}
\floatstyle{boxed}
\restylefloat{figure}


% Some colors:
\definecolor{dkcyan}{rgb}{0.1, 0.3, 0.3}
\definecolor{dkgreen}{rgb}{0,0.3,0}
\definecolor{olive}{rgb}{0.5, 0.5, 0.0}
\definecolor{dkblue}{rgb}{0,0.1,0.5}

\newcommand{\wrong}[1]{{\color{red} #1}}

\iftrue  % iftrue to reenable
\newcommand{\simon}[1]{{\color{dkblue} [{SLPJ}: #1]}}
\newcommand{\lennart}[1]{{\color{dkcyan} [{LA}: #1]}}
\else
\newcommand{\simon}[1]{}
\newcommand{\lennart}[1]{}
\fi

\newcommand{\vtt}[1]{{\tt #1}}

\begin{document}

%% Title information
\title{Verse}

                                        %% [Short Title] is optional;
                                        %% when present, will be used in
                                        %% header instead of Full Title.
% \titlenote{with title note}             %% \titlenote is optional;
                                        %% can be repeated if necessary;
                                        %% contents suppressed with 'anonymous'
% \subtitle{Subtitle}                     %% \subtitle is optional
% \subtitlenote{with subtitle note}       %% \subtitlenote is optional;
                                        %% can be repeated if necessary;
                                        %% contents suppressed with 'anonymous'


%% Author information
%% Contents and number of authors suppressed with 'anonymous'.
%% Each author should be introduced by \author, followed by
%% \authornote (optional), \orcid (optional), \affiliation, and
%% \email.
%% An author may have multiple affiliations and/or emails; repeat the
%% appropriate command.
%% Many elements are not rendered, but should be provided for metadata
%% extraction tools.

%% Author with single affiliation.
\author{Lennart Augustsson}
% \authornote{with author1 note}          %% \authornote is optional;
%                                         %% can be repeated if necessary
\orcid{nnnn-nnnn-nnnn-nnnn}             %% \orcid is optional
\affiliation{
  \position{Position1}
  \department{Epic Games}              %% \department is recommended
  \institution{Institution1}            %% \institution is required
  \streetaddress{Street1 Address1}
  \city{City1}
  \state{State1}
  \postcode{Post-Code1}
  \country{Country1}                    %% \country is recommended
}

\author{Simon Peyton-Jones}
% \authornote{with author1 note}          %% \authornote is optional;
%                                         %% can be repeated if necessary
\orcid{nnnn-nnnn-nnnn-nnnn}             %% \orcid is optional
\affiliation{
  \position{Position1}
  \institution{Microsoft Research}            %% \institution is required
  \streetaddress{Street1 Address1}
  \city{City1}
  \state{State1}
  \postcode{Post-Code1}
  \country{Country1}                    %% \country is recommended
}
\email{simonpj@@microsoft.com}          %% \email is recommended




%% Abstract
%% Note: \begin{abstract}...\end{abstract} environment must come
%% before \maketitle command
\begin{abstract}
  The abstract
\end{abstract}


%% 2012 ACM Computing Classification System (CSS) concepts
%% Generate at 'http://dl.acm.org/ccs/ccs.cfm'.
\begin{CCSXML}
<ccs2012>
<concept>
<concept_id>10011007.10011006.10011008</concept_id>
<concept_desc>Software and its engineering~General programming languages</concept_desc>
<concept_significance>500</concept_significance>
</concept>
<concept>
<concept_id>10003456.10003457.10003521.10003525</concept_id>
<concept_desc>Social and professional topics~History of programming languages</concept_desc>
<concept_significance>300</concept_significance>
</concept>
</ccs2012>
\end{CCSXML}

\ccsdesc[500]{Software and its engineering~General programming languages}
\ccsdesc[300]{Social and professional topics~History of programming languages}
%% End of generated code


%% Keywords
%% comma separated list
% \keywords{keyword1, keyword2, keyword3}  %% \keywords are mandatory in final camera-ready submission


%% \maketitle
%% Note: \maketitle command must come after title commands, author
%% commands, abstract environment, Computing Classification System
%% environment and commands, and keywords command.
\maketitle

%% Include unification (equality) among the primitives.
\newif\ifunify
\unifyfalse % change to unifytrue to include unification

\section{Introduction}

\lennart{Should we make all examples be valid actual Verse?  Possibly adding def{x}in.}

\section{A short introduction to Verse}

Verse is a statically typed functional-logic language: the main syntactic construct is the \emph{expression}.

\subsection{Expressions yield multiple values}

Verse's most salient and unusual feature is that an expression yields a (possibly-empty) \emph{sequence of values}.
\lennart{Sequences is a semantic concept; not something at the language level.}
For example:
\begin{itemize}
\item The expression \vtt{2} yields (a sequences consisting of) exactly one value, namely $2$.
\item The expression \vtt{(2 | 7 | 2)} yields a sequence of three values, namely $2$, $7$, and $2$.  We write this sequence $[2,7,2]$.
In general, the expression \vtt{($e_1$ | $e_2$)} yields all the values of $e_1$ followed by all the values of $e_2$.
\item The expression \vtt{:false} yields no values at all, the empty sequence $[]$.
\end{itemize}

An expression that yields no values is said to \emph{fail}; it if returns one or more, it is said to \emph{succeed}.

The semicolon operator propagates failure.
Specifically, the expression \vtt{($e_1$ ; $e_2$)}
fails if $e_1$ fails; and otherwise yields $e_2$. \simon{Is this right?  For example, consider
\vtt{((1 | 2) ; (3 | 4))}.  Does that yield $[3,4]$, or does it yield $[3,4,3,4]$?  The desugaring in Fig 2 suggests the latter...}

\subsection{Function application is lifted over the argment sequences}
\label{sec:lifting}

All functions are lifted to work over sequences, by taking the Cartesian product of their argument sequences.
For example, the expression \vtt{($e_1$ + $e_2$)} returns the sequence of all values $v_1 + v_2$ where
$v_1$ is drawn from the sequence yielded by $e_1$, and
$v_2$ is drawn from the sequence yielded by $e_2$.
Thus \vtt{(2 | 7) + (4 | 11)} evaluates to the sequence $[6, 13, 11, 18]$.

From this lifting semantics, you can see that failure propagates
upward; for example, \vtt{(:fail + $e$)} fails for any $e$, since there
are no values $v_1$ returned by \vtt{fail}.

\subsection{Comparisons and conditionals} \label{sec:bool}

In Verse, there are no Boolean values \vtt{true} and \vtt{false}.  Instead, in a conditional \vtt{if $e_1$ then $e_2$ else $e_3$},
if $e_1$ fails (yields no values), the conditional returns $e_3$.  If $e_1$ returns one or more values, the conditional returns $e_2$.
\cite{Icon}

The comparison operators such as $(>)$, $(<)$, $(=)$ and so on, use the same idiom: \vtt{($e_1$ < $e_2$)} returns $e_1$ if the condition
holds, and fails if it does not.  That makes \vtt{(if (3 > 4) then 10 else 20)} return $20$ as you would expect.

Amusingly, this semantics allows us to write
\vtt{(0 < x < 10)}, an expression that succeeds (returning $10$) if and only if \vtt{x} is between 0 and 10.
This only works because \vtt{(<)} is right associative, so the expression means \vtt{0 < (x<10)}. (Exercise for the reader:
what does \vtt{(0<x) < 10)} do?

What if $e_1$ and/or $e_2$ return a sequence of values?  Remembering the lifting semantics of \Cref{sec:lifting},
we take all pairs of values $(v_1,v_2)$ yielded by $e_1$ and $e_2$.  For each such pair, if $v_1 < v_2$ return $v_2$,
otherwise return nothing.   So \vtt{(1 | 8) < (2 | 10)} yeilds the sequence $[1,1,8]$.

\subsection{Variables}

In functional programming, a variable is bound to a specified immutable value as soon as it comes into scope.
Not so in Verse.  Instead, a variable can be brought into scope \emph{without a specified value}.  Here is an example:
\begin{verbatim}
  def x in x + (x = 7; 5)
\end{verbatim}
The construct \vtt{def x in $e$} brings \vtt{x} into scope in $e$.  Its semantics is unusual: the expression
\vtt{def x in $e$} binds $x$ successively to \emph{all possible values}, and returns all the values yielded
by $e$ under those bindings.  For which value(s) of \vtt{x} will the expression succeed?  Answer: because of the
\vtt{x=7} in the second argument of the \vtt{+}, all execuctions will fail except the one that binds \vtt{x}
to $7$.  So the expression yields just one value, namely 12.

Of course, no implementation could possibly work by binding \vtt{x} to every possible value and
evaluating $e$ -- there are an infinite set of such values.  But that is the semantics.
Notice that there is nothing special about \vtt{(=)}; it is just the ordinary equality operator
described in \Cref{sec:bool}.

It is easy to write an expression that will give any implementation a heart attack;
for example \vtt{def x in x} returns a sequence of all possible values.

In Verse, although an expression yields multiple values, a variable is always bound to a single value.

\begin{figure}

\begin{tabular}{c}
\begin{tabular}{|l|l|}
  \hline
  Notation &  \\
  \hline
  $x$ & variable \\
  $e,f,t$ & expression \\
  $i$ & natural number literal \\
  \hline
\end{tabular}

\\ \\[2mm]

\begin{tabular}{|l|l|}
  \hline
  Expression & Comment \\
  \hline
  $x$ & use variable \\
  $i$ & literal \\
\ifunify
  \vtt{$e_1$ = $e_2$} & unification \\
\fi
  \vtt{$e_1$[$e_2$]} & function application, may fail \\
  \vtt{$e_1$($e_2$)} & function application, may not fail \\
  \vtt{$x$ => $e$} & lambda expression \\
  \vtt{$e_1$ | $e_2$} & alternative \\
  \vtt{$e_1$,$\cdots$,$e_n$} & tuple formation, $n >= 2$, also $ n = 0$ in some syntactic contexts \\
  \vtt{if ($e_1$) then $e_2$ else $e_3$} & conditional \\
  \vtt{for ($e_1$) $e_2$} & array formation \\
  \hline
  \vtt{def\{$x_1$,...,$x_n$\} in $e$} & define variables in a scope, {\it only after extrusion} \\
  \hline
\end{tabular}
\end{tabular}

\caption{Core expressions}
\end{figure}

\begin{figure}
\begin{tabular}{c}
Constructs that are just syntactic sugar \\
\begin{tabular}{|l|l|l|}
  \hline
  Syntactic sugar (macro) & Expansion & Comment \\
  \hline
  \vtt{$x$ : $t$} & \vtt{$x$ := :$t$} & \\
  \vtt{typedef\{$e$\}} & \vtt{(x:any) => x = $e$} & \vtt{x} fresh \\
  \vtt{$e_1$ where $e_2$} & \vtt{($e_1$, $e_2$)[0]} & \\
  \vtt{case($e$) of \{$e_1$;$\cdots$;$e_n$\}} & \vtt{\{x:=$e$; $e_1$[x] || $\cdots$ || $e_n$[x]\}} & \vtt{x} fresh \\
  \vtt{\{ $e_0$;$\cdots$;$e_i$ \}} & \vtt{($e_0$,$\cdots$,$e_i$)[$i$]} & \\
  \hline \hline
  \vtt{$e_1$ \&\& $e_2$} & \vtt{if ($e_1$) then $e_2$ else :false} & {\it maybe not the right scope rules} \\
  \vtt{$e_1$ \&\& $e_2$} & \vtt{\{$e_1$; $e_2$\}} & {\it an alternate definition} \\
  \vtt{$e_1$ || $e_2$} & \vtt{if (x := $e_1$) then x else $e_2$} & \vtt{x} fresh \\
  \hline \hline
  \vtt{$e_1$:$t$ := $e_2$} & \vtt{$e_1$ := $t$[$e_2$]} & \\
  \vtt{($e_0$,$\cdots$,$e_n$) := $e$} & \vtt{\{x:=$e$; $e_0$:=x[0]; $\cdots$; $e_n$:=x[n]\}} & \vtt{x} fresh \\
  \vtt{$f$($e_1$) := $e_2$} & \vtt{$f$ := x => \{ x = $e_1$; $e_2$ \}} & \vtt{x} fresh \\
  \hline \hline
  \vtt{$e_1 \oplus e_2$} & \vtt{operator'$\oplus$'($e_1$,$e_2$) } & operator that cannot fail, e.g., \vtt{+}, \vtt{*} \\
  \vtt{$e_1 \oslash e_2$} & \vtt{operator'$\oslash$'[$e_1$,$e_2$] } & operator that can fail, e.g.,
\ifunify  \else
    \vtt{=},
\fi
    \vtt{/}, \vtt{<} \\
  \hline
\end{tabular}
\\ \\
Constructs that are eliminated by extrusion \\
\begin{tabular}{|l|l|}
  \hline
  Primitive & Comment \\
  \hline
  \vtt{$x$ := $e$} & define a variable \\
  \vtt{:}$t$ & range of a function \\
  \vtt{let ($e_1$) in $e_2$} & local bindings \\
  \vtt{do $e$} & limit scope \\
  \hline
\end{tabular}
\end{tabular}
\caption{Syntactic sugar}
\end{figure}

\begin{table}
\begin{tabular}{|l|l|}
  \hline
  Predefined name & Comment \\
  \hline
  \vtt{int} & identity on the integers \\
  \vtt{type} & identity on types \\
  \vtt{any} & identity function \\
  \vtt{array} & array type \\
  \vtt{tuple} & tuple type \\
  \vtt{arrow} & function type \\
  \vtt{wrong} & semantic error \\
  ... & ... \\
  \vtt{false} & \vtt{= ()} \\
  \vtt{void} & \vtt{= (:any) => false} \\
  \hline
\end{tabular}
\caption{Predefined identifiers}
\end{table}

\newcommand{\defdefin}[2]{\vtt{#1} $\rightsquigarrow$ \vtt{#2}}
\newcommand{\defarrow}[3]{\vtt{#1} $\leadsto$ \vtt{#2} ; #3}
\newcommand{\defempty}{$\varnothing$}
\newcommand{\defset}[1]{\{#1\}}
\newcommand{\defunion}[2]{$#1 \uplus #2$}
\newcommand{\defunions}[1]{$\biguplus #1$}
\newcommand{\defunique}[1]{$\# #1$}
\newcommand{\defafter}{\\ \vspace*{1em}}

\vspace{3em}

Figure \ref{extrusion-rules} has rules for collecting variable definitions to the beginning of blocks.
Notation: \defunion{}{} represents multiset union, $\#D$ is a predicate for multiset uniqueness,
\defempty is the empty multiset.
The relation \defarrow{$e$}{$e'$}{$D$} denotes that the topmost scope of $e$ contains the
\vtt{def}s for the variables in $D$, and $e'$ is the expression $e$ with all \vtt{def}s removed.
The relation \defdefin{$e$}{$e'$} denotes that $e'$ is equivalent to $e$,
but with all \vtt{def}s replaced by \vtt{def-in}.

\begin{figure}

\AxiomC{~}
\UnaryInfC{\defarrow{$x$}{$x$}{\defempty}}
\DisplayProof \defafter

\AxiomC{~}
\UnaryInfC{\defarrow{$i$}{$i$}{\defempty}}
\DisplayProof \defafter

\AxiomC{\defarrow{$e$}{$e'$}{$D$}}
\UnaryInfC{\defarrow{$x$ := $e$}{$x$ = $e'$}{\defunion{D}{\defset{x}}}}
\DisplayProof \defafter

\AxiomC{\defarrow{$t$}{$t'$}{$D$}}
\UnaryInfC{\defarrow{:$t$}{$t'$[def\{x\}in x]}{$D$}}
\DisplayProof \defafter

\ifunify
\AxiomC{\defarrow{$e_1$}{$e'_1$}{$D_1$}}
\AxiomC{\defarrow{$e_2$}{$e'_2$}{$D_2$}}
\BinaryInfC{\defarrow{$e_1$=$e_2$}{$e'_1$=$e'_2$}{\defunion{D_1}{D_2}}}
\DisplayProof \defafter
\fi

\AxiomC{\defarrow{$e_1$}{$e'_1$}{$D_1$}}
\AxiomC{\defarrow{$e_2$}{$e'_2$}{$D_2$}}
\BinaryInfC{\defarrow{$e_1$[$e_2$]}{$e'_1$[$e'_2$]}{\defunion{D_1}{D_2}}}
\DisplayProof \defafter

\AxiomC{\defarrow{$e_1$}{$e'_1$}{$D_1$}}
\AxiomC{\defarrow{$e_2$}{$e'_2$}{$D_2$}}
\BinaryInfC{\defarrow{$e_1$($e_2$)}{$e'_1$($e'_2$)}{\defunion{D_1}{D_2}}}
\DisplayProof \defafter

\AxiomC{\defdefin{$e$}{$e'$}}
\UnaryInfC{\defarrow{$x$=>$e$}{$x$=>$e'$}{\defempty}}
\DisplayProof \defafter

\AxiomC{\defdefin{$e_1$}{$e'_1$}}
\AxiomC{\defdefin{$e_2$}{$e'_2$}}
\BinaryInfC{\defarrow{$e_1$|$e_2$}{$e'_1$|$e'_2$}{\defempty}}
\DisplayProof \defafter

\AxiomC{\defarrow{$e_1$}{$e'_1$}{$D_1$}}
\AxiomC{$\cdots$}
\AxiomC{\defarrow{$e_n$}{$e'_n$}{$D_n$}}
\TrinaryInfC{\defarrow{$e_1$,$\cdots$,$e_n$}{$e'_1$,$\cdots$,$e'_n$}{\defunions{D_i}}}
\DisplayProof \defafter

\AxiomC{\defdefin{$e_1$}{$e'_1$}}
\AxiomC{\defdefin{$e_2$}{$e'_2$}}
\AxiomC{\defdefin{$e_3$}{$e'_3$}}
\TrinaryInfC{\defarrow{if($e_1$) then $e_2$ else $e_3$}{if($e'_1$) then $e'_2$ else $e'_3$}{\defempty}}
\DisplayProof
\\ The variables defined in the condition are also in scope in the \vtt{then} branch.
\defafter

\AxiomC{\defdefin{$e_1$}{$e'_1$}}
\AxiomC{\defdefin{$e_2$}{$e'_2$}}
\BinaryInfC{\defarrow{for($e_1$) $e_2$}{for($e'_1$) $e'_2$}{\defempty}}
\DisplayProof
\\ The variables defined in the domain are also in scope in body.
\defafter

\AxiomC{\defarrow{$e_1$}{$e'_1$}{$D_1$}}
\AxiomC{\defarrow{$e_2$}{$e'_2$}{$D_2$}}
\AxiomC{\defunique{D_1}}
\TrinaryInfC{\defarrow{let ($e_1$) in $e_2$}{def\{$D_1$\} in \{$e'_1$; $e'_2$\}}{$D_2$}}
\DisplayProof \defafter

\AxiomC{\defdefin{$e$}{$e'$}}
\UnaryInfC{\defarrow{do $e$}{$e'$}{\defempty}}
\DisplayProof \defafter

\AxiomC{\defarrow{$e_1$}{$e'_1$}{$D_1$}}
\AxiomC{$\cdots$}
\AxiomC{\defarrow{$e_n$}{$e'_n$}{$D_n$}}
\TrinaryInfC{\defarrow{\{$e_1$;$\cdots$;$e_n$\}}{\{$e'_1$;$\cdots$;$e'_n$\}}{\defunions{D_i}}}
\DisplayProof \defafter

\AxiomC{\defarrow{$e$}{$e'$}{$D$}}
\AxiomC{\defunique{D}}
\BinaryInfC{\defdefin{$e$}{def\{$D$\} in $e'$}}
\DisplayProof \defafter

\caption{Scope extrusion rules}
\label{extrusion-rules}
\end{figure}

\begin{figure}

\newcommand{\scinf}[2]{#1 $\vdash$ #2}

\AxiomC{$x \in \Gamma$}
\UnaryInfC{\scinf{$\Gamma$}{$x$}}
\DisplayProof \defafter

\AxiomC{}
\UnaryInfC{\scinf{$\Gamma$}{$i$}}
\DisplayProof \defafter

\AxiomC{\scinf{$\Gamma \cup D$}{\vtt{$e$}}}
\UnaryInfC{\scinf{$\Gamma$}{\vtt{def\{$D$\}in $e$}}}
\DisplayProof \defafter

\ifunify
\AxiomC{\scinf{$\Gamma$}{$e_1$}}
\AxiomC{\scinf{$\Gamma$}{$e_2$}}
\BinaryInfC{\scinf{$\Gamma$}{\vtt{$e_1$ = $e_2$}}}
\DisplayProof \defafter
\fi

\AxiomC{\scinf{$\Gamma$}{$e_1$}}
\AxiomC{\scinf{$\Gamma$}{$e_2$}}
\BinaryInfC{\scinf{$\Gamma$}{\vtt{$e_1$[$e_2$]}}}
\DisplayProof \defafter

\AxiomC{\scinf{$\Gamma$}{$e_1$}}
\AxiomC{\scinf{$\Gamma$}{$e_2$}}
\BinaryInfC{\scinf{$\Gamma$}{\vtt{$e_1$($e_2$)}}}
\DisplayProof \defafter

\AxiomC{\scinf{$\Gamma \cup \{x\}$}{\vtt{$e$}}}
\UnaryInfC{\scinf{$\Gamma$}{\vtt{$x$ => $e$}}}
\DisplayProof \defafter

\AxiomC{\scinf{$\Gamma$}{$e_1$}}
\AxiomC{\scinf{$\Gamma$}{$e_2$}}
\BinaryInfC{\scinf{$\Gamma$}{\vtt{$e_1$ | $e_2$}}}
\DisplayProof \defafter

\AxiomC{\scinf{$\Gamma$}{$e_1$}}
\AxiomC{$\cdots$}
\AxiomC{\scinf{$\Gamma$}{$e_n$}}
\TrinaryInfC{\scinf{$\Gamma$}{\vtt{$e_1$, $\cdots$, $e_n$}}}
\DisplayProof \defafter

\AxiomC{\scinf{$\Gamma \cup D$}{$e_1$}}
\AxiomC{\scinf{$\Gamma \cup D$}{$e_2$}}
\AxiomC{\scinf{$\Gamma$}{$e_3$}}
\TrinaryInfC{\scinf{$\Gamma$}{\vtt{if (def\{$D$\}in $e_1$) then $e_2$ else $e_3$}}}
\DisplayProof \defafter

\AxiomC{\scinf{$\Gamma \cup D$}{$e_1$}}
\AxiomC{\scinf{$\Gamma \cup D$}{$e_2$}}
\BinaryInfC{\scinf{$\Gamma$}{\vtt{for (def\{$D$\}in $e_1$) $e_2$}}}
\DisplayProof \defafter

\AxiomC{\scinf{$\Gamma$}{$e_1$}}
\AxiomC{$\cdots$}
\AxiomC{\scinf{$\Gamma$}{$e_n$}}
\TrinaryInfC{\scinf{$\Gamma$}{\vtt{\{$e_1$; $\cdots$; $e_n$\}}}}
\DisplayProof \defafter

\caption{Scoping rules}
\label{scoping-rules}
\end{figure}

\newcommand{\tuple}[1]{\langle #1 \rangle}
\newcommand{\Value}{V}
\newcommand{\Env}{E}

\begin{figure}
  $\Value = \mathbb{Z} + \tuple{\Value} + (\Value \rightarrow \Value^{*}) + \mathrm{WRONG}$ \\
  $\Env = Ident \rightarrow \Value$ \\
  where $\tuple{\Value}$ is a tuple (array) of values, $\Value^{*}$ is a (semantic) sequence (monad) of values.
\caption{Semantic domains}
\label{semantic-domains}
\end{figure}

\newcommand{\equa}[1]{\begin{equation*}#1\end{equation*}}
\newcommand{\mcompr}[2]{[ #1 \:|\: #2 ]}
\newcommand{\mbind}[2]{ #1 \leftarrow #2 }
\newcommand{\malt}{\diamond}
\newcommand{\munit}[1]{[#1]}
\newcommand{\mempty}{\munit{}}

\newcommand{\semapp}[2]{\mathcal{E}\llbracket #1 \rrbracket #2}
%\newcommand{\semeqn}[2]{\equa{ \semapp{#1}{\rho} = #2 }}
\newcommand{\semeqn}[2]{\semapp{#1}{\rho} & = #2 \\ }
\newcommand{\envapp}[2]{\mathcal{R}\llbracket #1 \rrbracket #2}
\newcommand{\liftrel}[1]{\hat{#1}}

\begin{figure}
  \begin{align*}
    \tuple{v_1,\cdots,v_n} && \text{create a tuple,}\, n >= 0 \\
    \munit{v_1,\cdots,v_n} && \text{create a sequence,}\, n >= 0 \\
    s_1 \malt s_2          && \text{sequence concatenation} \\
    \mcompr{e}{\mbind{x}{s}, \cdots} && \text{sequence comprehension} \\
    t_i                    && \text{index into a tuple} \\
    \mathrm{size}(t)       && \text{number of elements in a tuple}
  \end{align*}
\caption{Semantic notation}
\label{semantic-notation}
\end{figure}

\begin{figure}

  \begin{align*}
  \semapp{\cdot}{} &: \Env \rightarrow \Value^{*} \\
  \semeqn{x}{\munit{\rho(x)}}
  \semeqn{i}{\munit{i}}
\ifunify
  \semeqn{e_1 \vtt{=} e_2}{\mathrm{lift}(\hat{=}, \semapp{e_1}{\rho}, \semapp{e_2}{\rho})}
\fi
  \semeqn{e_1[e_2]}{\mathrm{lift}(\mathrm{apply}, \semapp{e_1}{\rho}, \semapp{e_2}{\rho})}
  \semeqn{e_1(e_2)}{\mathrm{lift}(\mathrm{call}, \semapp{e_1}{\rho}, \semapp{e_2}{\rho})}
  \semeqn{x \vtt{=>} e}{\munit{\lambda v . \semapp{e} (\rho \circ (x \mapsto v))}}
  \semeqn{e_1 \vtt{|} e_2}{\semapp{e_1}{\rho} \malt \semapp{e_2}{\rho}}
  \semeqn{e_1,\cdots,e_n}{
    \mcompr{\tuple{v_1,\cdots,v_n}}
          {\mbind{v_1}{\semapp{e_1}{\rho}}, \cdots, \mbind{v_n}{\semapp{e_n}{\rho}}}
          }
  \semeqn{\vtt{if(}e_1\vtt{)\:then}\: e_2 \: \vtt{else}\: e_3}{
      \begin{cases*}
        \semapp{e_2}{\rho'} & if $\envapp{e_1}{s\, \rho} = \munit{\rho', \cdots}$ \\
        \semapp{e_3}{\rho} & otherwise
      \end{cases*}
      }
  \semeqn{\vtt{for(}e_1\vtt{)}e_2}{\mathrm{tuples}(\mcompr{\semapp{e_2}{\rho'}}{\mbind{\rho'}{\envapp{e_1}{\rho}}})}
  \semeqn{\vtt{def\{}x_1,\cdots,x_n\vtt{\}in}\,e}
         {\mcompr{\semapp{e}{(\rho \circ \{ x_1 \mapsto v_1, \cdots, x_n \mapsto v_n\})}} {\mbind{v_1}{univ},\cdots,\mbind{v_n}{univ}}}
  \end{align*}

  \equa { \envapp{\cdot}{} : \Env \rightarrow \Env^{*} }
  \equa {
    \begin{aligned}
      \envapp{e}{\rho} =
      \rho \circ  \{ & x_1 \mapsto \semapp{\vtt{def\{}x_1,\cdots,x_n\vtt{\} in (}e',x_1\vtt{)[1]}}{\rho}, \\
                     & \cdots, \\
                     & x_n \mapsto \semapp{\vtt{def\{}x_1,\cdots,x_n\vtt{\} in (}e',x_n\vtt{)[1]}}{\rho}, \\
      & \mathrm{where}\, e = \vtt{def\{}x_1,\cdots,x_n\vtt{\} in}\,e'
    \end{aligned}
    }
  
  \equa{\mathrm{lift}(f, w_1, w_2) = \mcompr{v}{\mbind{v_1}{w_1}, \mbind{v_2}{w_2}, \mbind{v}{f(v_1, v_2)}}}

\ifunify
  \equa{\hat{f}(x,y) =
    \begin{cases*}
      \munit{x} & if $f(x,y)$ \\
      \mempty & otherwise
    \end{cases*}
    }
\fi

  \equa{\mathrm{apply}(f,x) = 
      \begin{cases*}
        f(x) & if $f \in \Value \rightarrow \Value^{*}$ \\
        \munit{f_x} & if $f \in \tuple{\Value}, x \in \mathbb{Z}, x \geqslant 0, x < \mathrm{size}(f)$ \\
        \mempty & otherwise, if $f \in \tuple{\Value}$ \\
        \munit{\mathrm{WRONG}} & otherwise
      \end{cases*}
  }
  \equa{\mathrm{call}(f,x) = 
      \begin{cases*}
        apply(f,x) & if $apply(f,x) \ne \mempty$ \\
        \munit{\mathrm{WRONG}} & otherwise
      \end{cases*}
  }

  \begin{multline*}
    \mathrm{tuples}(\munit{
      \munit{v_{11},v_{12},\cdots,v_{1n_1}},
      \munit{v_{21},v_{22},\cdots,v_{2n_2}},
      \cdots,
      \munit{v_{m1},v_{m2},\cdots,v_{mn_m}}}) = \\
    \munit{
      \tuple{v_{11},v_{21},\cdots,v_{m1}},
      \tuple{v_{11},v_{21},\cdots,v_{m2}},
      \cdots,
      \tuple{v_{1n_1},\cdots,v_{mn_m}}}
  \end{multline*}

  \lennart{This is all slighly wrong, since it doesn't propagate WRONG correctly.
    Perhaps * should be a monad that can throw an exception when something is WRONG.
  }

\caption{Expression semantics}
\label{expression-semantics}
\end{figure}

\newcommand{\checkcase}[3]{\begin{cases*} #1 & if $#2$ \\ #3 & otherwise \end{cases*}}
\newcommand{\isfunction}[1]{#1 \in \Value \rightarrow \Value^{*}}

\begin{table}
  \begin{tabular}{|l|l|}
    \hline
    Name & Semantic value \\
    \hline
    \vtt{any} & $x \mapsto \munit{x}$ \\
    \vtt{false} & $x \mapsto \mempty$ \\
    \vtt{int} & $x \mapsto \checkcase{\munit{x}}{x \in \mathbb{Z}}{\mempty} $ \\
    \vtt{tuple} & $f \mapsto x \mapsto
      \begin{cases*}
        \mathrm{tuples}(f_1(x_1),\cdots,f_n(x_n)) & if $f = \tuple{f_1,\cdots,f_n}, x = \tuple{x_1,\cdots,x_n}$ \\
          & $ \isfunction{f_1}, \cdots, \isfunction{f_n} $ \\
          \mempty & otherwise
      \end{cases*}
      $ \\
    \vtt{array} & $f \mapsto x \mapsto
      \checkcase{\mathrm{tuples}(f(x_1),\cdots,f(x_n))}
                {x = \tuple{x_1,\cdots,x_n}, \isfunction{f}}
                {\mempty} $ \\
    \vtt{arrow} & $ t \mapsto u \mapsto f \mapsto x \mapsto $ \\
         & $ \quad \checkcase{\mcompr{v}{\mbind{y}{t(x)}, \mbind{z}{f(y)}, \mbind{v}{u(z)}}}{\isfunction{t}, \isfunction{u}, \isfunction{f}}{\mempty}
                  $ \\
    \vtt{operator'+'} & $x \mapsto \checkcase{\munit{x_1+x_2}}{x = \tuple{x_1,x_2}, x_1 \in \mathbb{Z}, x_2 \in \mathbb{Z}}{\mempty}$ \\
    \vtt{operator'/'} & $x \mapsto \checkcase{\munit{\frac{x_1}{x_2}}}{x = \tuple{x_1,x_2}, x_1 \in \mathbb{Z}, x_2 \in \mathbb{Z}, x_2 \ne 0}{\mempty}$ \\
\ifunify
\else
    \vtt{operator'='}  & $x \mapsto \checkcase{\munit{x_1}}{x = \tuple{x_1,x_2}, x_1 = x_2}{\mempty}$ \\
\fi
    \vtt{operator'<'}  & $x \mapsto \checkcase{\munit{x_1}}{x = \tuple{x_1,x_2}, x_1 < x_2}{\mempty}$ \\
    \hline
  \end{tabular}
\caption{Some predefined identifiers.  Note that several of these could be defined in Verse}
\label{predefined}
\end{table}

\begin{acks}                            %% acks environment is optional
                                        %% contents suppressed with 'anonymous'
\end{acks}

%% Bibliography
\bibliography{bibliography}


\end{document}

